\documentclass[11pt]{article}
\usepackage{amsmath,amssymb,amsthm,hyperref}
\newtheorem{theorem}{Theorem}
\newtheorem{proposition}{Proposition}
\newtheorem{corollary}{Corollary}
\newcommand{\R}{\mathbb R}
\newcommand{\E}{\mathbb E}
\newcommand{\PP}{\mathbb P}
\newcommand{\one}{\mathbf 1}
\title{A flat-outlier obstruction to sparse partition certificates}
\author{Sparse partitioning research session}
\date{November 15, 2026}
\begin{document}
\maketitle

\paragraph{Status.}
This is an elementary, completely proved obstruction, not a claim of a new
spectral technique or of historical priority. Its mechanism is the classical
motif-doubling/flat-band construction. Kawamoto explicitly applies motif
doubling to nonbacktracking partitioning [K, Sections 4--6]. The Cartesian
product flat-band generator of Sabri--Youssef [SY, Theorem 3.1] is a closely
related primary antecedent. The purpose here is to make the certification
limitation precise even when both individual eigenvectors and their full
eigenspace have the strongest possible usual incoherence, and to give an
exact statistical ambiguity bound. This is not a counterexample to a
stochastic block model or an independent random-lift theorem.

\paragraph{Construction.}
Let $R$ be a connected simple $r$-regular graph on $k$ vertices, with
$r\geq1$, and $F$ a connected simple $q$-regular graph on $h$ vertices,
with $q\geq1$. Define a graph $G=G(R,F)$
on vertices $(i,s,a)$, where $i\in V(R)$, $s\in\{+,-\}$, $a\in V(F)$.
Inside each $(i,s)$ put a copy of $F$. For every edge $ij$ of $R$ and every
$a$, put all four edges between $\{(i,+,a),(i,-,a)\}$ and
$\{(j,+,a),(j,-,a)\}$. Then $G$ is simple, connected, and
$d=q+2r$ regular, with $N=2hk$ vertices. If $R$ and $F$ are bipartite, so
is $G$: add their two bipartite colors modulo two.

Let $J_2$ be the all-ones matrix of order two. In the displayed coordinates,
\begin{equation}
 A_G=I_k\otimes I_2\otimes A_F+A_R\otimes J_2\otimes I_h.
 \label{tensor}
\end{equation}
For each $i$ define the unit vector
\[
 e_i(j,s,a)=\frac{\one_{i=j}\,s}{\sqrt{2h}},
 \qquad s\in\{+1,-1\}.
\]
Write $E=\operatorname{span}\{e_i:1\leq i\leq k\}$ and let $P_E$ be its
orthogonal projector. For a rank-$k$ subspace of $\R^N$ define
$\mu(E)=(N/k)\max_v(P_E)_{vv}$. Always $\mu(E)\geq1$.

\begin{theorem}[Exact outliers with optimal incoherence]
For the construction above:
\begin{enumerate}
\item $A_G|_E=qI_E$. Every $e_i$ has support $2h$.
\item $(P_E)_{vv}=k/N=1/(2h)$ for every vertex, so $\mu(E)=1$ exactly.
\item If a real Hadamard matrix of order $k$ exists, in particular if
      $k=2^m$, $E$ has an orthonormal eigenbasis $f_1,\ldots,f_k$ with
      $|f_j(v)|=N^{-1/2}$ for every $j,v$. Consequently each such vector
      has inverse participation ratio exactly $1/N$, its minimum possible
      value among unit vectors in $\R^N$.
\item If $q>2\sqrt{d-1}$, these are adjacency outliers above the spectrum
      of the universal cover, the $d$-regular tree. For
      $H(t)=(t^2-1)I+D-tA_G$, they have negative Bethe--Hessian eigenvalue
      $2(d-1)-q\sqrt{d-1}$ at $t=\sqrt{d-1}$.
\item Suppose $F=K_{q,q}$, $q>2r$, and $A_R$ is invertible. Then $E$ is
      the entire $q$-eigenspace of $A_G$, not just a subspace of it.
\end{enumerate}
\end{theorem}
\begin{proof}
The two terms in \eqref{tensor} commute. The $+$/$-$ antisymmetric direction
annihilates $J_2$, while the constant eigenvector of $F$ has eigenvalue $q$.
This proves the first assertion. Disjoint supports and normalization give
$(P_E)_{vv}=1/(2h)$. If $U$ is a Hadamard matrix divided by $\sqrt{k}$,
then $f_j=\sum_iU_{ij}e_i$ gives the stated orthonormal basis. Its entries
have modulus $1/\sqrt{2hk}$, proving its participation ratio. On a regular
graph $H(\sqrt{d-1})=2(d-1)I-\sqrt{d-1}A_G$, proving the fourth assertion.

More generally the full adjacency spectrum, including multiplicity, is
\begin{equation}
\begin{gathered}
 \{\lambda_j(F):1\leq j\leq h\}\quad\text{each repeated $k$ times},\\
 \{\lambda_j(F)+2\theta_i(R):1\leq j\leq h,\ 1\leq i\leq k\}.
\end{gathered}
 \label{spec}
\end{equation}
For $F=K_{q,q}$ its eigenvalues are $q,-q,0$. In the first collection only
its simple top eigenvalue gives $q$, yielding exactly $E$. In the second,
$q+2\theta=q$ would require $\theta=0$, ruled out by invertibility;
$-q+2\theta=q$ requires $\theta=q>r$; and $2\theta=q$ is impossible since
$q>2r$. Thus the full $q$-eigenspace is $E$.
\end{proof}

\begin{proposition}[Uniform expansion is compatible with the obstruction]
Let $\gamma_R=r-\lambda_2(A_R)>0$ and
$\gamma_F=q-\lambda_2(A_F)>0$. Then the ordinary Laplacian gap is exactly
\[
 \lambda_2(L_G)=\min\{2r,2\gamma_R,\gamma_F\}.
\]
In particular fixed-degree expander families $R$ and a fixed connected $F$
give a fixed-degree expander family $G$.
\end{proposition}
\begin{proof}
In the antisymmetric collection in \eqref{spec}, the largest eigenvalue is
$q=d-2r$. In the symmetric collection the unique top value is
$q+2r=d$; the second largest value is
$\max\{q+2\lambda_2(A_R),\lambda_2(A_F)+2r\}$.
Taking $d$ minus the largest nonprincipal value proves the formula.
\end{proof}

\begin{proposition}[Exponential equitable ambiguity and exact Bayes bound]
For every $\sigma\in\{\pm1\}^k$, the labels
\[
 x_\sigma(i,s,a)=s\sigma_i
\]
form a balanced bisection. Every vertex has exactly $q+r$ neighbors of its
own label and $r$ of the other label, and $A_Gx_\sigma=qx_\sigma$.
Choose $\sigma_i$ independently uniformly, and optionally apply a uniformly
random permutation $\pi$ to all vertices and labels. In the latter version
the marginal planted labeling is uniform over all balanced labelings.
From the observed
unlabeled graph, no estimator $\widehat x\in\{\pm1\}^N$ has expected
sign-invariant overlap greater than
\[
 \frac{\E|\sigma_1+\cdots+\sigma_k|}{k}\leq\frac1{\sqrt{k}}.
\]
The same bound holds even if $\pi$, the cells, and their $+$/$-$ orientation
are revealed. In this stronger oracle problem the first expression is the
exact optimum. For even $k$ it equals
\[
 2^{-k}\binom{k}{k/2}\sim\sqrt{\frac{2}{\pi k}}.
\]
For any estimator and any $\delta>0$,
\[
 \PP\left\{N^{-1}|\langle\widehat x,x_\sigma\rangle|\geq\delta\right\}
 \leq 2\exp(-k\delta^2/2).
\]
Thus weak recovery is impossible in this specified model, even when $q$
is strictly above the universal-cover spectral threshold.
\end{proposition}
\begin{proof}
Each cell has equally many labels of each sign. The $q$ neighbors inside
$F$ all have the same label as a vertex. For each of its $r$ neighboring
cells there is one same-label and one opposite-label neighbor, irrespective
of $\sigma$. This proves the degree assertion and eigenvalue equation.
There are $2^{k-1}$ distinct bisections modulo global interchange. A uniform
permutation of any fixed balanced labeling is uniform over all balanced
labelings, proving the marginal-prior assertion.

Condition on $\pi$ and the estimator's independent randomness, even revealing
$\pi$ to the estimator. The observed graph contains no dependence on
$\sigma$. Put
\[
 w_i=\sum_{a\in V(F)}\bigl[\widehat x(\pi(i,+,a))-
                   \widehat x(\pi(i,-,a))\bigr],\qquad |w_i|\leq2h.
\]
Then $\langle\widehat x,\pi x_\sigma\rangle=\sum_iw_i\sigma_i$.
The function $\E|\sum_iw_i\sigma_i|$ is convex in each $w_i$, so its maximum
on the box $[-2h,2h]^k$ occurs at a vertex. At all vertices it equals
$2h\E|\sum_i\sigma_i|$. Taking canonical labels with $w_i=2h$ attains
this maximum in the oracle problem. Divide by $N=2hk$. Cauchy--Schwarz
gives the stated upper bound; the central binomial formula for even $k$
follows by summing the binomial distribution (or the adjacent-binomial
telescoping identity). Hoeffding's inequality and
$\sum_iw_i^2\leq4h^2k$ give the probability bound. These arguments also
apply to randomized estimators after averaging the conditioning.
\end{proof}

\paragraph{The same degree constraints and conditional first moments.}
In the randomly permuted version, conditional on any balanced planted
labeling, every vertex has within-label degree $a=q+r$ and across-label
degree $b=r$. Permutation invariance within each label class therefore gives
the exact one-edge marginals $a/(N/2-1)$ within a class and $b/(N/2)$ across
classes. The severe edge correlations, not a weak first-moment signal,
cause the impossibility. Independence or another identifying stochastic
assumption cannot be replaced by these degree constraints alone.

\paragraph{An explicit fixed-degree family.}
The following device avoids silently assuming the expander $R$ is
nonsingular. Start with any family of connected simple cubic expanders
$S_m$ on $2^m$ vertices, with $3-\lambda_2(S_m)\geq\gamma$ for some
$0<\gamma\leq1$.
Let $B_m=S_m[K_2]$ be the lexicographic product: replace every vertex by
an edge and every base edge by $K_{2,2}$. It has degree seven and spectrum
$2\theta+1$ for $\theta\in\sigma(S_m)$, together with $-1$ repeated
$|S_m|$ times. It is nonsingular: otherwise $-1/2$ would be an eigenvalue
of an integer matrix, impossible since a rational algebraic integer is an
integer. Its smallest eigenvalue is at least $-5$, and its nonprincipal
largest eigenvalue is at most $7-2\gamma$.

Take the bipartite double cover $R_m$ of $B_m$. It is simple, seven-regular,
nonsingular, and connected because $B_m$ is connected and contains a
triangle. Its adjacency eigenvalues are the positives and negatives of
those of $B_m$, so
$\gamma_{R_m}\geq\min\{2\gamma,2\}>0$.
It has $k=2^{m+2}$ vertices. Now take $F=K_{15,15}$. Then $G(R_m,F)$ is a
simple connected bipartite $29$-regular expander, $N=60k$, whose full
$15$-eigenspace has rank $k$, optimal incoherence one, and a perfectly flat
real orthonormal eigenbasis, while
\[
 15>2\sqrt{28},\qquad
 H(\sqrt{28})|_E=(56-15\sqrt{28})I_E<0.
\]
Classical random regular graph results supply the cubic expander family;
no new existence result for expanders is needed here.

\paragraph{A small exact diagnostic.}
Take $R$ to be the three-dimensional cube, $r=3$, and $F=K_{7,7}$.
Then $k=8,h=14,d=13,N=224$, $A_R$ is nonsingular, and the full
$7$-eigenspace has rank eight. Its eigenvalue exceeds $2\sqrt{12}$,
the Laplacian gap is exactly four, and the corresponding large
nonbacktracking root is exactly four since $4+12/4=7$.

\paragraph{Nonbacktracking eigenvectors.}
Use the directed-edge matrix
\[
 (Bz)(u,v)=\sum_{w\sim v,w\ne u}z(v,w).
\]
Assume $q>2\sqrt{d-1}$. If $A_Gf=qf$ and $\beta$ is the larger root of
$\beta+(d-1)/\beta=q$, so $\beta>\sqrt{d-1}>1$, then
\[
 z_f(u,v)=f(v)-\beta^{-1}f(u),\qquad Bz_f=\beta z_f,
 \quad \|z_f\|^2=(d-2)(1-\beta^{-2})\|f\|^2.
\]
Thus the flat orthonormal vertex basis maps, after one common
normalization, to an orthonormal basis of right nonbacktracking outliers
whose infinity norms are $O(N^{-1/2})$ for fixed $d$.
For a flat $f$, their participation ratio is exactly
\[
 \frac{(q+r)(1-\beta^{-1})^4+r(1+\beta^{-1})^4}
 {N(d-2)^2(1-\beta^{-2})^2}.
\]
For the $d=13$ example this is $179/(1815N)$, and the infinity norm is
$5/\sqrt{165N}$, with $13N$ directed edges.

\begin{corollary}[The same obstruction in independent random lifts]
Let $G_m=G(R_m,K_{15,15})$ be the explicit $29$-regular bipartite base
family above, with $k_m=|R_m|=2^{m+2}$ and $D_m=60k_m$. Choose any lift
degree $n_m$ with $k_m/n_m\to0$. Sample an independent uniform $n_m$-lift
$\mathcal G_m$ of $G_m$. Its total size is $M_m=60k_m n_m$.
Choose uniform independent hidden signs $\sigma\in\{\pm1\}^{k_m}$,
constant on each lifted cell as above, and optionally hide all coordinates
by a uniform vertex permutation. Then:
\begin{enumerate}
\item Every hidden labeling is a balanced exact $(22,7)$ equitable
      bisection, and the optimal expected sign-invariant overlap is at most
      $\E|\sum_i\sigma_i|/k_m\sim\sqrt{2/(\pi k_m)}\to0$.
      Equality holds in the stronger oracle problem revealing all covering
      coordinates and the final vertex permutation.
\item With probability $1-O(k_m/n_m)$, the \emph{entire} $15$-eigenspace
      of $A_{\mathcal G_m}$ is the lifted old space $E_m$, of rank $k_m$.
      Its projector has diagonal $k_m/M_m=1/(60n_m)$ everywhere, its
      incoherence is exactly one, and it has a real eigenbasis with every
      entry of modulus $M_m^{-1/2}$.
\item On the same event the graph is a connected simple bipartite expander
      with a uniform positive ordinary Laplacian gap. Moreover its new
      adjacency spectrum is asymptotically almost Ramanujan in probability.
\item For each fixed radius $\ell$, the proportion of vertices whose
      radius-$\ell$ neighborhood contains a cycle tends to zero in
      probability. Thus the rooted local limit is the $29$-regular tree.
\end{enumerate}
The rank tends to infinity, but can do so arbitrarily slowly relative to
total graph size; in particular $\operatorname{rank}(E_m)/M_m\to0$.
\end{corollary}
\begin{proof}
The only probabilistic spectral input is the existing uniform bound [GL]:
for every deterministic base of size $D$ and maximum degree $d$,
an independent $n$-lift satisfies
\[
 \PP\{\|A_{\rm new}\|\geq\rho(T)+\varepsilon\}
 \leq C\frac Dn\left(\frac d\varepsilon\right)^9.
\]
Here $d=29$, $\rho(T)=2\sqrt{28}$, and
$\varepsilon_0=(15-2\sqrt{28})/2>0$ is a fixed constant. This gives an event
of probability at least $1-C'k_m/n_m$ on which
$\|A_{\rm new}\|<2\sqrt{28}+\varepsilon_0<15$.
The old adjacency spectrum is exactly the base spectrum. Since the full
base $15$-eigenspace is $E$, on this event the full lifted $15$-eigenspace
is precisely its old lift. A normalized old vector is its base value divided
by $\sqrt{n_m}$ on every sheet, proving all flatness and projector claims.

Every covering preserves the exact within-/across-label neighbor counts
of a fiber-constant labeling. The graph is sampled without any dependence
on $\sigma$. Conditional on the random graph and all oracle coordinates,
the previous Bayes proof applies with cell half-size $30n_m$ in place of
$h=30$, proving the same optimal overlap and its concentration bound.

Let the base Laplacian gap be at least $c>0$, independent of $m$. On the
displayed event every old nonprincipal adjacency eigenvalue is at most
$29-c$, and every new one is less than $15$. Hence
$\lambda_2(L_{\mathcal G_m})\geq\min\{c,14\}>0$. In particular the graph
is connected: an extra connected component would give another eigenvector
of adjacency eigenvalue $29$. Simplicity and bipartiteness hold for every
lift of this simple bipartite base. For each fixed $\varepsilon>0$, the
same input bound tends to zero because $D_m/n_m\to0$, establishing the
asymptotic almost-Ramanujan statement on the new space.

Finally explore a fixed-radius ball at a chosen lifted vertex by exposing
its matching entries. At most a constant $C_{29,\ell}$ entries are exposed.
At each exposure the unmatched endpoint is uniform among at least
$n_m-C_{29,\ell}$ remaining vertices in its base fiber. The probability of
collision with any previously exposed vertex is therefore $O(1/n_m)$,
uniformly in the base vertex and base size. A union bound gives an
$O(1/n_m)$ collision probability for the whole ball. Averaging over roots
and applying Markov's inequality proves the final assertion.
\end{proof}

\paragraph{Scope of the obstruction and a real remaining problem.}
For the fixed-motif deterministic family the outlier rank is linear in $N$;
for the lifted family it grows while its fraction tends to zero. The result does not
contradict fixed-rank spectral clustering theory, an independent-edge SBM,
a uniform regular block model, or independent random-lift delocalization.
In particular the rank-aware incoherence bound
$P_{vv}\leq\mu\operatorname{rank}(P)/N$ does not say $P_{vv}=O(1/N)$
when the rank grows. The construction refutes a certificate based only on
an outlier, a flat displayed eigenbasis, bounded degree, global expansion,
rank-aware incoherence, or even local tree-likeness and independent lifting
when the old base has many ambiguous fine types. It does not refute certificates which also
control outlier rank or inspect local decomposition of the full projector.

A substantial remaining target is a quantitative, basis-invariant
decomposition theorem for independently lifted or semirandom sparse block
models: split the informative low-rank subspace from finite-defect modes
using the off-diagonal/local structure of the spectral projector, with
recovery valid after near-degenerate defect modes mix. This note proves no
such algorithm. The recent fixed-base random-lift local law [HX] and clean
SBM Bethe--Hessian analysis [SZ] must be included in any priority comparison.

\begin{thebibliography}{9}
\bibitem{K} T. Kawamoto, \emph{Localized eigenvectors of the non-backtracking
matrix}, J. Stat. Mech. (2016), 023404.
\url{https://arxiv.org/abs/1505.07543v3}.
\bibitem{SY} M. Sabri and P. Youssef, \emph{Flat bands of periodic graphs},
\url{https://arxiv.org/abs/2304.06465v2}.
\bibitem{HX} J. Huang and R. Xu, \emph{Local Law and Eigenvector
Delocalization for Random Lifts}, submitted October 27, 2026.
\url{https://arxiv.org/abs/2610.37463v1}.
\bibitem{SZ} L. Stephan and Y. Zhu, \emph{Community detection with the
Bethe-Hessian}, v3, November 4, 2026.
\url{https://arxiv.org/abs/2411.02835v3}.
\bibitem{GL} \emph{A sharp worst-case base-size threshold for random cubic
lifts}, repository research note, November 15, 2026, first theorem
(which applies to general bounded-degree bases), and the companion proof
\texttt{generator-free-tail-proof.md} in its \texttt{prior/} directory.
This is a refinement of the rational-trace method of Chen, Garza-Vargas,
Tropp and van Handel, \emph{A new approach to strong convergence},
\url{https://arxiv.org/abs/2405.16026}.
\end{thebibliography}
\end{document}
